\section{Interfaces entre visión y lenguaje: de Q-Former a la proyección lineal}

Al conectar image con un LLM, la primera pregunta no es «¿el modelo entendió?», sino cómo entran las características continuas del encoder visual en el token space del language model. BLIP-2 y LLaVA ofrecen dos formas típicas de hacer ese puente: uno usa queries aprendibles y un Q-Former para comprimir y alinear; el otro proyecta linealmente y de forma directa los token visuales.

\subsection{BLIP-2: consultar el espacio visual congelado con un Q-Former}

BLIP-2 coloca un Q-Former entre un frozen image encoder y un frozen LLM. Sean las características visuales $V\in\mathbb{R}^{m\times d_v}$ y las queries aprendibles $Q_0\in\mathbb{R}^{q\times d}$; la cross attention es
\[
Q'=\operatorname{softmax}\!\left(\frac{(Q_0W_Q)(VW_K)^\top}{\sqrt d}\right)VW_V.
\]
Cuando $q\ll m$, el Q-Former funciona a la vez como cuello de botella de información y como alineación entre espacios; después envía al LLM un número reducido de query outputs.

\lecturefigure{slide-13-blip2.jpg}{BLIP-2: el Q-Former sirve de puente entre un frozen image encoder y un frozen LLM}{Grabación oficial de Stanford Online 00:36:18}

\begin{knowledgebox}{El Q-Former resuelve dos problemas distintos}
El primero es la \textbf{compression}: extraer un número fijo de queries a partir de una gran cantidad de image patches. El segundo es la \textbf{alignment}: llevar el feature space de tipo CLIP al embedding space que el LLM puede consumir. No se trata de que el LLM vuelva a aprender por sí mismo todo el backbone visual.
\end{knowledgebox}

\teachervoice{00:35:00--00:36:55, el ponente describe BLIP-2 como pionero en conectar de forma eficaz CLIP con un LLM; el énfasis está en que «entre dos espacios congelados hace falta un puente entrenable».}

\subsection{LLaVA: por qué una projection simple se volvió un baseline sólido}

LLaVA usa un vision encoder $f_v$ y una projection $W_p$:
\[
H_v=f_v(I),\qquad Z_v=H_vW_p,
\]
y luego concatena $Z_v$ con los text embeddings como entrada del LLM. Reduce los módulos adicionales y facilita el instruction tuning end-to-end, por lo que pronto se convirtió en un baseline habitual de los VLM.

\lecturefigure{slide-14-llava.jpg}{LLaVA: alinear las características visuales con el embedding de lenguaje mediante una projection}{Grabación oficial de Stanford Online 00:36:55}

\begin{warningbox}{Que la linear projection sea simple no significa que no pierda información}
La projection solo garantiza que las dimensiones coincidan. La resolución del encoder visual, el patch size, el pretraining objective, la projector capacity y los instruction data determinan qué detalles llegan al LLM; el texto pequeño, el OCR, el grounding y los GUI icons suelen ser los primeros en revelar el cuello de botella.
\end{warningbox}

\subsection{Taxonomy de interfaces: compresión, alineación, conservación y alta resolución}

Antes de pasar a la serie Cog, todos los VLM bridge pueden compararse con cuatro preguntas: ¿cuántos token visuales se conservan? ¿Comparten parámetros la visión y el lenguaje? ¿Se daña la capacidad de lenguaje original? ¿Dónde se paga el cost de la entrada de alta resolución?

\begin{knowledgebox}{Tabla comparativa de VLM bridge}
\begin{tabularx}{\textwidth}{L{0.16\textwidth} L{0.22\textwidth} L{0.23\textwidth} X}
\toprule
Método & Interfaz visual & Ruta de parámetros & Compromiso principal \\
\midrule
BLIP-2 & Q-Former queries & bridge entrenable, backbones congelables & compresión fuerte, pero el query bottleneck es evidente \\
LLaVA & linear/MLP projector & los token visuales entran en la LLM path & línea base simple y sólida; costo alto en alta resolución \\
CogVLM & vision experts & los image/text token pasan en parte por parámetros distintos & intenta conservar el comportamiento de lenguaje; más parámetros \\
CogAgent & high-res cross attention & desvío ligero de alta resolución & orientado a GUI/OCR; sistema más complejo \\
Vary & múltiples vocabularios visuales/feature ensemble & fusión de características de varios encoder & OCR más fuerte; aumentan los costos de interfaz y entrenamiento \\
GLM-4V slide & stride convolution & mixed training tras la compresión & estructura simple; los resultados dependen de los datos y del protocol \\
\bottomrule
\end{tabularx}
\end{knowledgebox}

\subsection{Resumen de la sección}

En esencia, la interfaz multimodal es una cuestión de information bottleneck y parameter sharing. El Q-Former consulta y comprime de forma explícita; LLaVA proyecta directamente. Los modelos posteriores se diversifican en torno a «conservar la capacidad de lenguaje» y «procesar los detalles de alta resolución».
